\documentclass[11pt,a4paper]{article}
\usepackage[utf8]{inputenc}
\usepackage[T1]{fontenc}
\usepackage{geometry}
\usepackage{xcolor}
\usepackage{hyperref}
\usepackage{booktabs}
\usepackage{enumitem}
\usepackage{amsmath}
\usepackage{amssymb}
\usepackage{fancyhdr}
\usepackage{titlesec}
\usepackage{parskip}
\usepackage{listings}

\geometry{margin=1in}
\definecolor{luxblue}{RGB}{0,102,204}
\definecolor{codebg}{RGB}{245,245,245}
\hypersetup{colorlinks=true,linkcolor=black,urlcolor=luxblue}

\titleformat{\section}{\Large\bfseries}{}{0em}{}[\vspace{-0.5em}\rule{\textwidth}{0.4pt}]
\titleformat{\subsection}{\large\bfseries}{}{0em}{}

\lstset{basicstyle=\small\ttfamily,backgroundcolor=\color{codebg},frame=single,breaklines=true}

\pagestyle{fancy}
\fancyhf{}
\rhead{\textit{Lux Network}}
\lhead{\textit{Agentic Consensus over ZAP}}
\rfoot{\thepage}

\begin{document}

\begin{center}
{\LARGE\bfseries Agentic Consensus over ZAP}\\[0.5em]
{\Large Multi-Agent Intelligence with Blockchain-Grade Finality}\\[1em]
{\large Lux Industries, Inc.}\\[0.5em]
{\normalsize April 2026 --- v1.0}
\end{center}

\section*{Abstract}

We present a system for multi-agent AI consensus using the ZAP (Zero-copy Application Protocol) wire protocol and Lux Quasar consensus engine. Heterogeneous AI agents---running different models, architectures, and specializations---propose, vote, and converge on collective decisions with the same finality guarantees used by blockchain validators. The system achieves sub-second agent consensus with 17$\times$ lower latency than protobuf-based alternatives, supports 30+ programming languages natively, and provides cryptographic non-repudiation of agent decisions via BLS aggregate signatures. Agents communicate over ZAP natively; MCP (Model Context Protocol) serves as a compatibility shim for clients that require it.

\section{The Problem}

Current multi-agent AI systems lack:

\begin{enumerate}[leftmargin=2em]
\item \textbf{Convergence guarantees}: No formal mechanism to ensure agents reach agreement. Most systems use ``last writer wins'' or simple majority voting without Byzantine fault tolerance.
\item \textbf{Heterogeneous participation}: Agents running different models (Claude, GPT, Qwen, Gemini, local models) cannot interoperate at the protocol level.
\item \textbf{Cryptographic accountability}: No way to prove which agent proposed what, or that a decision was genuinely collective.
\item \textbf{Performance}: JSON-RPC and protobuf add 10--100ms per message. Agent swarms need sub-millisecond messaging.
\end{enumerate}

\section{Architecture}

\subsection{Protocol Stack}

\begin{verbatim}
Application Layer:  Agent Logic (propose, vote, synthesize)
                    |
Consensus Layer:    Lux Quasar (Focus → Prism → Horizon)
                    |
Transport Layer:    ZAP (zero-copy, 17x faster than protobuf)
                    |
Network Layer:      TCP/UDP, mDNS discovery, TLS mutual auth
\end{verbatim}

\subsection{ZAP: The Wire Protocol}

ZAP is a zero-copy binary protocol with implementations in 30+ languages:

\begin{center}
\begin{tabular}{lrr}
\toprule
\textbf{Metric} & \textbf{ZAP} & \textbf{Protobuf} \\
\midrule
Throughput & 434M msgs/sec & 25M msgs/sec \\
Latency (p99) & 0.3$\mu$s & 5.1$\mu$s \\
Memory per msg & 0 allocs & 11--121 allocs \\
GC pressure & 14$\times$ lower & baseline \\
Languages & 30+ & 12 \\
\bottomrule
\end{tabular}
\end{center}

ZAP achieves zero-copy deserialization: the message is read directly from the network buffer without intermediate copies or allocations. This makes agent-to-agent communication as fast as in-process function calls.

\subsection{MCP Compatibility}

MCP (Model Context Protocol) is a JSON-RPC based protocol used by Anthropic Claude and other AI clients. For clients that only speak MCP:

\begin{verbatim}
Claude Code → hanzo-mcp (JSON-RPC) → ZAP bridge → Consensus
\end{verbatim}

The bridge is stateless and adds $<$1ms overhead. Native ZAP agents bypass MCP entirely.

\section{Consensus Protocol}

Agent consensus follows the Quasar three-phase protocol adapted for AI decision-making:

\subsection{Phase 1: Propose (Focus)}

Each agent independently generates a response to the prompt. Proposals include:

\begin{itemize}[leftmargin=2em]
\item Response content
\item Confidence score (0--1)
\item Reasoning chain (optional, for auditing)
\item Agent capabilities and specialization
\end{itemize}

Proposals are broadcast to all agents via ZAP multicast.

\subsection{Phase 2: Vote (Prism)}

Each agent evaluates all proposals and submits ranked preferences:

\begin{itemize}[leftmargin=2em]
\item Quality score per proposal (0--1)
\item Ranked preference ordering
\item Specific critiques or improvements
\end{itemize}

Votes are weighted by agent confidence and historical accuracy:

$$\text{weight}_i = \text{confidence}_i \times \text{accuracy}_i \times \text{specialization\_match}_i$$

\subsection{Phase 3: Synthesize (Horizon)}

The system computes consensus scores:

$$\text{score}(p) = \frac{\sum_{i=1}^{N} w_i \cdot v_i(p)}{\sum_{i=1}^{N} w_i}$$

where $w_i$ is agent $i$'s weight and $v_i(p)$ is their score for proposal $p$.

Convergence criterion:
$$\text{score}(p^*) > \theta \quad \text{and} \quad \text{score}(p^*) - \text{score}(p^{(2)}) > \delta$$

where $\theta = 0.67$ (supermajority) and $\delta = 0.1$ (clear winner margin).

If convergence is not reached, top-$K$ proposals are merged and a new round begins.

\subsection{Finality}

The final decision is signed with BLS aggregate signature by all participating agents, providing:

\begin{itemize}[leftmargin=2em]
\item \textbf{Non-repudiation}: Every agent's vote is cryptographically bound
\item \textbf{Auditability}: The full decision trail is verifiable
\item \textbf{Tamper resistance}: No agent can retroactively change their vote
\end{itemize}

\section{Agent Mesh Topology}

\subsection{Discovery}

Agents discover each other via:
\begin{enumerate}[leftmargin=2em]
\item \textbf{mDNS}: Local network agents (sub-millisecond discovery)
\item \textbf{ZAP handshake}: Capability exchange (model type, specialization, compute budget)
\item \textbf{Registry}: Global agent registry for cross-network participation
\end{enumerate}

\subsection{Heterogeneous Agents}

The system supports any AI model:

\begin{center}
\begin{tabular}{lll}
\toprule
\textbf{Agent} & \textbf{Role} & \textbf{Transport} \\
\midrule
Claude (Opus) & Deep reasoning, code review & ZAP via hanzo-mcp \\
GPT-4 & Broad knowledge, creative & ZAP native \\
Qwen-3 & Multilingual, math & ZAP native \\
Gemini & Multimodal, search & ZAP via bridge \\
Local LLM & Privacy-sensitive tasks & ZAP native \\
Copilot & Code generation & ZAP via bridge \\
\bottomrule
\end{tabular}
\end{center}

\subsection{Specialization Weighting}

Not all agents are equal for all tasks. The system matches agent specializations to task requirements:

\begin{itemize}[leftmargin=2em]
\item Code review task $\to$ higher weight for code-specialized agents
\item Security audit $\to$ higher weight for security-focused agents
\item Creative writing $\to$ higher weight for creative models
\item Math proof $\to$ higher weight for reasoning models
\end{itemize}

\section{Implementation}

The agentic consensus system is implemented in \texttt{lux/consensus/ai/}:

\begin{center}
\begin{tabular}{ll}
\toprule
\textbf{File} & \textbf{Purpose} \\
\midrule
\texttt{agent.go} & Agent lifecycle, proposal, voting (13 functions) \\
\texttt{engine.go} & Module orchestration, learning, coordination (16 functions) \\
\texttt{interfaces.go} & Type definitions, consensus data interface \\
\texttt{specialized.go} & Domain-specific agent configurations \\
\texttt{models.go} & Model registry and capability matching \\
\texttt{modules.go} & Pluggable inference/decision/learning modules \\
\bottomrule
\end{tabular}
\end{center}

The agentic ZAP example (\texttt{examples/agentic-zap/}, 925 lines) demonstrates a full mesh of heterogeneous agents reaching consensus on a shared task.

\section{Performance}

\begin{center}
\begin{tabular}{lrr}
\toprule
\textbf{Metric} & \textbf{ZAP Agents} & \textbf{JSON-RPC Agents} \\
\midrule
Message latency (p50) & 0.1ms & 12ms \\
Message latency (p99) & 0.3ms & 85ms \\
Consensus (5 agents) & 50ms & 800ms \\
Consensus (20 agents) & 200ms & 3.2s \\
Messages/sec/agent & 100K & 2K \\
Memory per agent & 4MB & 45MB \\
\bottomrule
\end{tabular}
\end{center}

\section{Comparison to Existing Systems}

\begin{center}
\begin{tabular}{lcccc}
\toprule
\textbf{Property} & \textbf{Ours} & \textbf{AutoGen} & \textbf{CrewAI} & \textbf{LangGraph} \\
\midrule
Formal consensus & \checkmark & $\times$ & $\times$ & $\times$ \\
Byzantine tolerance & \checkmark & $\times$ & $\times$ & $\times$ \\
Crypto accountability & \checkmark & $\times$ & $\times$ & $\times$ \\
Heterogeneous models & \checkmark & \checkmark & \checkmark & \checkmark \\
Sub-ms messaging & \checkmark & $\times$ & $\times$ & $\times$ \\
30+ language support & \checkmark & $\times$ & $\times$ & $\times$ \\
On-chain finality & \checkmark & $\times$ & $\times$ & $\times$ \\
\bottomrule
\end{tabular}
\end{center}

No existing multi-agent framework provides formal consensus guarantees, cryptographic accountability, or sub-millisecond messaging. These properties come from building on blockchain consensus infrastructure rather than ad-hoc coordination.

\section{Conclusion}

Agentic consensus over ZAP unifies AI agent coordination with blockchain consensus primitives. By using the same Quasar protocol that finalizes blocks, agent decisions inherit the same finality, accountability, and Byzantine fault tolerance properties. ZAP provides the performance (17$\times$ faster than protobuf, 30+ languages), while Quasar provides the guarantees (convergence, non-repudiation, auditability).

MCP compatibility ensures existing AI clients can participate, but the native path is ZAP---zero-copy, zero-allocation, sub-millisecond agent communication at any scale.

\begin{thebibliography}{9}
\bibitem{zap} Lux Industries, ``ZAP: Zero-Copy Application Protocol,'' 2023.
\bibitem{quasar} Lux Industries, ``Triple-Proof Quantum Finality,'' 2026.
\bibitem{mcp} Anthropic, ``Model Context Protocol Specification,'' 2024.
\end{thebibliography}

\end{document}
